
\documentclass{article}
\usepackage{colortbl}
\usepackage{makecell}
\usepackage{multirow}
\usepackage{supertabular}

\begin{document}

\newcounter{utterance}

\centering \large Interaction Transcript for game `hot\_air\_balloon', experiment `air\_balloon\_survival\_de\_reasoning off\_hard', episode 3 with deepseek{-}chat{-}v3.1{-}t1.0.
\vspace{24pt}

{ \footnotesize  \setcounter{utterance}{1}
\setlength{\tabcolsep}{0pt}
\begin{supertabular}{c@{$\;$}|p{.15\linewidth}@{}p{.15\linewidth}p{.15\linewidth}p{.15\linewidth}p{.15\linewidth}p{.15\linewidth}}
    \# & $\;$A & \multicolumn{4}{c}{Game Master} & $\;\:$B\\
    \hline

    \theutterance \stepcounter{utterance}  
    & & \multicolumn{4}{p{0.6\linewidth}}{
        \cellcolor[rgb]{0.9,0.9,0.9}{
            \makecell[{{p{\linewidth}}}]{
                \texttt{\tiny{[P1$\langle$GM]}}
                \texttt{Du nimmst an einem kollaborativen Verhandlungspiel Teil.} \\
\\ 
\texttt{Zusammen mit einem anderen Teilnehmer musst du dich auf eine Reihe von Gegenständen entscheiden, die behalten werden. Jeder von euch hat eine persönliche Verteilung über die Wichtigkeit der einzelnen Gegenstände. Jeder von euch hat eine eigene Meinung darüber, wie wichtig jeder einzelne Gegenstand ist (Gegenstandswichtigkeit). Du kennst die Wichtigkeitsverteilung des anderen Spielers nicht. Zusätzlich siehst du, wie viel Aufwand jeder Gegenstand verursacht.  } \\
\texttt{Ihr dürft euch nur auf eine Reihe von Gegenständen einigen, wenn der Gesamtaufwand der ausgewählten Gegenstände den Maximalaufwand nicht überschreitet:} \\
\\ 
\texttt{Maximalaufwand = 3629} \\
\\ 
\texttt{Hier sind die einzelnen Aufwände der Gegenstände:} \\
\\ 
\texttt{Aufwand der Gegenstände = \{"C76": 5, "C38": 793, "C56": 21, "A03": 842, "A07": 313, "A83": 842, "C98": 631, "C08": 226, "C62": 87, "C00": 761, "C10": 230, "C81": 287, "B38": 698, "C03": 641, "C32": 882\}} \\
\\ 
\texttt{Hier ist deine persönliche Verteilung der Wichtigkeit der einzelnen Gegenstände:} \\
\\ 
\texttt{Werte der Gegenstandswichtigkeit = \{"C76": 138, "C38": 583, "C56": 868, "A03": 822, "A07": 783, "A83": 65, "C98": 262, "C08": 121, "C62": 508, "C00": 780, "C10": 461, "C81": 484, "B38": 668, "C03": 389, "C32": 808\}} \\
\\ 
\texttt{Ziel:} \\
\\ 
\texttt{Dein Ziel ist es, eine Reihe von Gegenständen auszuhandeln, die dir möglichst viel bringt (d. h. Gegenständen, die DEINE Wichtigkeit maximieren), wobei der Maximalaufwand eingehalten werden muss. Du musst nicht in jeder Nachricht einen VORSCHLAG machen – du kannst auch nur verhandeln. Alle Taktiken sind erlaubt!} \\
\\ 
\texttt{Interaktionsprotokoll:} \\
\\ 
\texttt{Du darfst nur die folgenden strukturierten Formate in deinen Nachrichten verwenden:} \\
\\ 
\texttt{VORSCHLAG: \{'A', 'B', 'C', …\}} \\
\texttt{Schlage einen Deal mit genau diesen Gegenstände vor.} \\
\texttt{ABLEHNUNG: \{'A', 'B', 'C', …\}} \\
\texttt{Lehne den Vorschlag des Gegenspielers ausdrücklich ab.} \\
\texttt{ARGUMENT: \{'...'\}} \\
\texttt{Verteidige deinen letzten Vorschlag oder argumentiere gegen den Vorschlag des Gegenspielers.} \\
\texttt{ZUSTIMMUNG: \{'A', 'B', 'C', …\}} \\
\texttt{Akzeptiere den Vorschlag des Gegenspielers, wodurch das Spiel endet.} \\
\\ 
\\ 
\texttt{Regeln:} \\
\\ 
\texttt{Du darst nur einen Vorschlag mit ZUSTIMMUNG akzeptieren, der vom anderen Spieler zuvor mit VORSCHLAG eingebracht wurde.} \\
\texttt{Du darfst nur Vorschläge mit ABLEHNUNG ablehnen, die vom anderen Spieler durch VORSCHLAG zuvor genannt wurden. } \\
\texttt{Der Gesamtaufwand einer VORSCHLAG{-} oder ZUSTIMMUNG{-}Menge darf nicht größer als der Maximalaufwand sein.  } \\
\texttt{Offenbare deine versteckte Wichtigkeitsverteilung nicht.} \\
\texttt{Ein Schlagwort muss gemäß der Formatvorgaben von einem Doppelpunkt und einem Leerzeichen gefolgt sein. Das Argument muss eine Python{-}Menge mit 0 oder mehr Strings sein.  } \\
\texttt{Die Form muss also SCHLAGWORT: \{...\} sein.} \\
\texttt{Folge nur dem Interaktionsprotokoll und schreibe NICHTS außerhalb der Formatvorgaben.} \\
\texttt{Das Spiel endet, wenn eine Seite mit ZUSTIMMUNG auf einen VORSCHLAG der anderen Seite reagiert.  } \\
\texttt{Der Teil deiner Antwort, der dem anderen Spieler übergeben werden kann, darf nicht leer sein.  } \\
\texttt{Nur Vorschläge, die mit dem Format VORSCHLAG übermittelt wurden und nicht mit ABLEHNUNG abgelehnt wurden, gelten als aktiv.  } \\
\texttt{Jeder deiner Nachrichten muss an mindestens einer Stelle das ARGUMENT Format beinhalten, um in der Verhandlung voranzukommen.} \\
            }
        }
    }
    & \\ \\

    \theutterance \stepcounter{utterance}  
    & \multicolumn{4}{p{0.6\linewidth}}{
        \cellcolor[rgb]{0.9,0.9,0.9}{
            \makecell[{{p{\linewidth}}}]{
                \texttt{\tiny{[P1$\rangle$GM]}}
                \texttt{ARGUMENT: \{'Ich möchte eine faire und für beide vorteilhafte Einigung erzielen. Lass uns Gegenstände priorisieren, die für uns jeweils einen hohen Wert bei moderaten Kosten bieten.'\}  } \\
\texttt{VORSCHLAG: \{'C56', 'A03', 'A07', 'C00', 'C10', 'C81', 'B38', 'C32', 'C76', 'C62'\}} \\
            }
        }
    }
    & & \\ \\

    \theutterance \stepcounter{utterance}  
    & & \multicolumn{4}{p{0.6\linewidth}}{
        \cellcolor[rgb]{0.9,0.9,0.9}{
            \makecell[{{p{\linewidth}}}]{
                \texttt{\tiny{[GM$\rangle$P2]}}
                \texttt{Du nimmst an einem kollaborativen Verhandlungspiel Teil.} \\
\\ 
\texttt{Zusammen mit einem anderen Teilnehmer musst du dich auf eine Reihe von Gegenständen entscheiden, die behalten werden. Jeder von euch hat eine persönliche Verteilung über die Wichtigkeit der einzelnen Gegenstände. Jeder von euch hat eine eigene Meinung darüber, wie wichtig jeder einzelne Gegenstand ist (Gegenstandswichtigkeit). Du kennst die Wichtigkeitsverteilung des anderen Spielers nicht. Zusätzlich siehst du, wie viel Aufwand jeder Gegenstand verursacht.  } \\
\texttt{Ihr dürft euch nur auf eine Reihe von Gegenständen einigen, wenn der Gesamtaufwand der ausgewählten Gegenstände den Maximalaufwand nicht überschreitet:} \\
\\ 
\texttt{Maximalaufwand = 3629} \\
\\ 
\texttt{Hier sind die einzelnen Aufwände der Gegenstände:} \\
\\ 
\texttt{Aufwand der Gegenstände = \{"C76": 5, "C38": 793, "C56": 21, "A03": 842, "A07": 313, "A83": 842, "C98": 631, "C08": 226, "C62": 87, "C00": 761, "C10": 230, "C81": 287, "B38": 698, "C03": 641, "C32": 882\}} \\
\\ 
\texttt{Hier ist deine persönliche Verteilung der Wichtigkeit der einzelnen Gegenstände:} \\
\\ 
\texttt{Werte der Gegenstandswichtigkeit = \{"C56": 65, "A03": 121, "C32": 138, "A07": 262, "C00": 389, "B38": 461, "C38": 484, "C62": 508, "C81": 583, "C10": 668, "C03": 780, "C98": 783, "C76": 808, "C08": 822, "A83": 868\}} \\
\\ 
\texttt{Ziel:} \\
\\ 
\texttt{Dein Ziel ist es, eine Reihe von Gegenständen auszuhandeln, die dir möglichst viel bringt (d. h. Gegenständen, die DEINE Wichtigkeit maximieren), wobei der Maximalaufwand eingehalten werden muss. Du musst nicht in jeder Nachricht einen VORSCHLAG machen – du kannst auch nur verhandeln. Alle Taktiken sind erlaubt!} \\
\\ 
\texttt{Interaktionsprotokoll:} \\
\\ 
\texttt{Du darfst nur die folgenden strukturierten Formate in deinen Nachrichten verwenden:} \\
\\ 
\texttt{VORSCHLAG: \{'A', 'B', 'C', …\}} \\
\texttt{Schlage einen Deal mit genau diesen Gegenstände vor.} \\
\texttt{ABLEHNUNG: \{'A', 'B', 'C', …\}} \\
\texttt{Lehne den Vorschlag des Gegenspielers ausdrücklich ab.} \\
\texttt{ARGUMENT: \{'...'\}} \\
\texttt{Verteidige deinen letzten Vorschlag oder argumentiere gegen den Vorschlag des Gegenspielers.} \\
\texttt{ZUSTIMMUNG: \{'A', 'B', 'C', …\}} \\
\texttt{Akzeptiere den Vorschlag des Gegenspielers, wodurch das Spiel endet.} \\
\\ 
\\ 
\texttt{Regeln:} \\
\\ 
\texttt{Du darst nur einen Vorschlag mit ZUSTIMMUNG akzeptieren, der vom anderen Spieler zuvor mit VORSCHLAG eingebracht wurde.} \\
\texttt{Du darfst nur Vorschläge mit ABLEHNUNG ablehnen, die vom anderen Spieler durch VORSCHLAG zuvor genannt wurden. } \\
\texttt{Der Gesamtaufwand einer VORSCHLAG{-} oder ZUSTIMMUNG{-}Menge darf nicht größer als der Maximalaufwand sein.  } \\
\texttt{Offenbare deine versteckte Wichtigkeitsverteilung nicht.} \\
\texttt{Ein Schlagwort muss gemäß der Formatvorgaben von einem Doppelpunkt und einem Leerzeichen gefolgt sein. Das Argument muss eine Python{-}Menge mit 0 oder mehr Strings sein.  } \\
\texttt{Die Form muss also SCHLAGWORT: \{...\} sein.} \\
\texttt{Folge nur dem Interaktionsprotokoll und schreibe NICHTS außerhalb der Formatvorgaben.} \\
\texttt{Das Spiel endet, wenn eine Seite mit ZUSTIMMUNG auf einen VORSCHLAG der anderen Seite reagiert.  } \\
\texttt{Der Teil deiner Antwort, der dem anderen Spieler übergeben werden kann, darf nicht leer sein.  } \\
\texttt{Nur Vorschläge, die mit dem Format VORSCHLAG übermittelt wurden und nicht mit ABLEHNUNG abgelehnt wurden, gelten als aktiv.  } \\
\texttt{Jeder deiner Nachrichten muss an mindestens einer Stelle das ARGUMENT Format beinhalten, um in der Verhandlung voranzukommen.} \\
\\ 
\\ 
\texttt{Du erhältst nun die erste Nachricht des anderen Spielers.} \\
\\ 
\texttt{ARGUMENT: \{'Ich möchte eine faire und für beide vorteilhafte Einigung erzielen. Lass uns Gegenstände priorisieren, die für uns jeweils einen hohen Wert bei moderaten Kosten bieten.'\}  } \\
\texttt{VORSCHLAG: \{'C56', 'A03', 'A07', 'C00', 'C10', 'C81', 'B38', 'C32', 'C76', 'C62'\}} \\
            }
        }
    }
    & \\ \\

    \theutterance \stepcounter{utterance}  
    & & & \multicolumn{4}{p{0.6\linewidth}}{
        \cellcolor[rgb]{0.9,0.9,0.9}{
            \makecell[{{p{\linewidth}}}]{
                \texttt{\tiny{[GM$\langle$P2]}}
                \texttt{ARGUMENT: \{'Ich schätze deinen Ansatz, eine faire Einigung zu finden. Allerdings enthält dein Vorschlag mehrere Gegenstände mit sehr hohem Aufwand, die den Budgetrahmen stark belasten.'\}  } \\
\texttt{VORSCHLAG: \{'C76', 'C56', 'A07', 'C62', 'C81', 'C10', 'C00', 'C03', 'C98'\}} \\
            }
        }
    }
     \\ \\

    \theutterance \stepcounter{utterance}  
    & & \multicolumn{4}{p{0.6\linewidth}}{
        \cellcolor[rgb]{0.9,0.9,0.9}{
            \makecell[{{p{\linewidth}}}]{
                \texttt{\tiny{[P1$\langle$GM]}}
                \texttt{ARGUMENT: \{'Ich schätze deinen Ansatz, eine faire Einigung zu finden. Allerdings enthält dein Vorschlag mehrere Gegenstände mit sehr hohem Aufwand, die den Budgetrahmen stark belasten.'\}  } \\
\texttt{VORSCHLAG: \{'C76', 'C56', 'A07', 'C62', 'C81', 'C10', 'C00', 'C03', 'C98'\}} \\
            }
        }
    }
    & \\ \\

    \theutterance \stepcounter{utterance}  
    & \multicolumn{4}{p{0.6\linewidth}}{
        \cellcolor[rgb]{0.9,0.9,0.9}{
            \makecell[{{p{\linewidth}}}]{
                \texttt{\tiny{[P1$\rangle$GM]}}
                \texttt{ARGUMENT: \{'Dein Vorschlag enthält einige wertvolle Gegenstände, aber C98 und C03 haben für mich eine relativ geringe Wichtigkeit im Vergleich zu ihrem Aufwand. Ich würde lieber B38 und C32 einbeziehen, die mir wichtiger sind.'\}  } \\
\texttt{VORSCHLAG: \{'C76', 'C56', 'A07', 'C62', 'C81', 'C10', 'C00', 'B38', 'C32'\}} \\
            }
        }
    }
    & & \\ \\

    \theutterance \stepcounter{utterance}  
    & & \multicolumn{4}{p{0.6\linewidth}}{
        \cellcolor[rgb]{0.9,0.9,0.9}{
            \makecell[{{p{\linewidth}}}]{
                \texttt{\tiny{[GM$\rangle$P2]}}
                \texttt{ARGUMENT: \{'Dein Vorschlag enthält einige wertvolle Gegenstände, aber C98 und C03 haben für mich eine relativ geringe Wichtigkeit im Vergleich zu ihrem Aufwand. Ich würde lieber B38 und C32 einbeziehen, die mir wichtiger sind.'\}  } \\
\texttt{VORSCHLAG: \{'C76', 'C56', 'A07', 'C62', 'C81', 'C10', 'C00', 'B38', 'C32'\}} \\
            }
        }
    }
    & \\ \\

    \theutterance \stepcounter{utterance}  
    & & & \multicolumn{4}{p{0.6\linewidth}}{
        \cellcolor[rgb]{0.9,0.9,0.9}{
            \makecell[{{p{\linewidth}}}]{
                \texttt{\tiny{[GM$\langle$P2]}}
                \texttt{ARGUMENT: \{'Ich verstehe deine Präferenz für B38 und C32. Allerdings haben diese Gegenstände einen sehr hohen Aufwand (698 und 882), was unser Budget stark belastet. Mein Vorschlag priorisiert Gegenstände mit besserem Wert{-}Aufwand{-}Verhältnis.'\}  } \\
\texttt{ABLEHNUNG: \{'C76', 'C56', 'A07', 'C62', 'C81', 'C10', 'C00', 'B38', 'C32'\}  } \\
\texttt{VORSCHLAG: \{'C76', 'C56', 'A07', 'C62', 'C81', 'C10', 'C00', 'C03', 'C98'\}} \\
            }
        }
    }
     \\ \\

    \theutterance \stepcounter{utterance}  
    & & \multicolumn{4}{p{0.6\linewidth}}{
        \cellcolor[rgb]{0.9,0.9,0.9}{
            \makecell[{{p{\linewidth}}}]{
                \texttt{\tiny{[P1$\langle$GM]}}
                \texttt{ARGUMENT: \{'Ich verstehe deine Präferenz für B38 und C32. Allerdings haben diese Gegenstände einen sehr hohen Aufwand (698 und 882), was unser Budget stark belastet. Mein Vorschlag priorisiert Gegenstände mit besserem Wert{-}Aufwand{-}Verhältnis.'\}  } \\
\texttt{ABLEHNUNG: \{'C76', 'C56', 'A07', 'C62', 'C81', 'C10', 'C00', 'B38', 'C32'\}  } \\
\texttt{VORSCHLAG: \{'C76', 'C56', 'A07', 'C62', 'C81', 'C10', 'C00', 'C03', 'C98'\}} \\
            }
        }
    }
    & \\ \\

    \theutterance \stepcounter{utterance}  
    & \multicolumn{4}{p{0.6\linewidth}}{
        \cellcolor[rgb]{0.9,0.9,0.9}{
            \makecell[{{p{\linewidth}}}]{
                \texttt{\tiny{[P1$\rangle$GM]}}
                \texttt{ARGUMENT: \{'Ich sehe den Wert in deinem Vorschlag, aber C03 und C98 bieten mir nicht genug Wert für ihren Aufwand. Mein Counter{-}Vorschlag behält die effizienten Gegenstände bei und fügt A03 hinzu, das einen ausgezeichneten Wert für mich bietet.'\}} \\
\texttt{VORSCHLAG: \{'C76', 'C56', 'A07', 'C62', 'C81', 'C10', 'C00', 'B38', 'A03'\}} \\
            }
        }
    }
    & & \\ \\

    \theutterance \stepcounter{utterance}  
    & & \multicolumn{4}{p{0.6\linewidth}}{
        \cellcolor[rgb]{0.9,0.9,0.9}{
            \makecell[{{p{\linewidth}}}]{
                \texttt{\tiny{[GM$\rangle$P2]}}
                \texttt{ARGUMENT: \{'Ich sehe den Wert in deinem Vorschlag, aber C03 und C98 bieten mir nicht genug Wert für ihren Aufwand. Mein Counter{-}Vorschlag behält die effizienten Gegenstände bei und fügt A03 hinzu, das einen ausgezeichneten Wert für mich bietet.'\}} \\
\texttt{VORSCHLAG: \{'C76', 'C56', 'A07', 'C62', 'C81', 'C10', 'C00', 'B38', 'A03'\}} \\
            }
        }
    }
    & \\ \\

    \theutterance \stepcounter{utterance}  
    & & & \multicolumn{4}{p{0.6\linewidth}}{
        \cellcolor[rgb]{0.9,0.9,0.9}{
            \makecell[{{p{\linewidth}}}]{
                \texttt{\tiny{[GM$\langle$P2]}}
                \texttt{ARGUMENT: \{'A03 hat einen sehr hohen Aufwand (842) für seinen Wert. B38 ist ebenfalls sehr teuer. Mein Vorschlag mit C03 und C98 bietet ein besseres Gesamtwert{-}Aufwand{-}Verhältnis bei gleichem Budget.'\}  } \\
\texttt{ABLEHNUNG: \{'C76', 'C56', 'A07', 'C62', 'C81', 'C10', 'C00', 'B38', 'A03'\}  } \\
\texttt{VORSCHLAG: \{'C76', 'C56', 'A07', 'C62', 'C81', 'C10', 'C00', 'C03', 'C98'\}} \\
            }
        }
    }
     \\ \\

    \theutterance \stepcounter{utterance}  
    & & \multicolumn{4}{p{0.6\linewidth}}{
        \cellcolor[rgb]{0.9,0.9,0.9}{
            \makecell[{{p{\linewidth}}}]{
                \texttt{\tiny{[P1$\langle$GM]}}
                \texttt{ARGUMENT: \{'A03 hat einen sehr hohen Aufwand (842) für seinen Wert. B38 ist ebenfalls sehr teuer. Mein Vorschlag mit C03 und C98 bietet ein besseres Gesamtwert{-}Aufwand{-}Verhältnis bei gleichem Budget.'\}  } \\
\texttt{ABLEHNUNG: \{'C76', 'C56', 'A07', 'C62', 'C81', 'C10', 'C00', 'B38', 'A03'\}  } \\
\texttt{VORSCHLAG: \{'C76', 'C56', 'A07', 'C62', 'C81', 'C10', 'C00', 'C03', 'C98'\}} \\
            }
        }
    }
    & \\ \\

    \theutterance \stepcounter{utterance}  
    & \multicolumn{4}{p{0.6\linewidth}}{
        \cellcolor[rgb]{0.9,0.9,0.9}{
            \makecell[{{p{\linewidth}}}]{
                \texttt{\tiny{[P1$\rangle$GM]}}
                \texttt{ARGUMENT: \{'I understand you're focused on cost efficiency, but C03 (389 value/641 cost) and C98 (262 value/631 cost) have poor value{-}to{-}cost ratios for me. I'm willing to compromise but need better value items.'\}} \\
\texttt{VORSCHLAG: \{'C76', 'C56', 'A07', 'C62', 'C81', 'C10', 'C00', 'B38', 'C08'\}} \\
            }
        }
    }
    & & \\ \\

    \theutterance \stepcounter{utterance}  
    & & \multicolumn{4}{p{0.6\linewidth}}{
        \cellcolor[rgb]{0.9,0.9,0.9}{
            \makecell[{{p{\linewidth}}}]{
                \texttt{\tiny{[GM$\rangle$P2]}}
                \texttt{ARGUMENT: \{'I understand you're focused on cost efficiency, but C03 (389 value/641 cost) and C98 (262 value/631 cost) have poor value{-}to{-}cost ratios for me. I'm willing to compromise but need better value items.'\}} \\
\texttt{VORSCHLAG: \{'C76', 'C56', 'A07', 'C62', 'C81', 'C10', 'C00', 'B38', 'C08'\}} \\
            }
        }
    }
    & \\ \\

    \theutterance \stepcounter{utterance}  
    & & & \multicolumn{4}{p{0.6\linewidth}}{
        \cellcolor[rgb]{0.9,0.9,0.9}{
            \makecell[{{p{\linewidth}}}]{
                \texttt{\tiny{[GM$\langle$P2]}}
                \texttt{ARGUMENT: \{'C08 hat einen moderaten Aufwand (226) und bietet mir einen guten Wert, was diesen Vorschlag akzeptabler macht als deine vorherigen. Allerdings ist B38 mit 698 Aufwand immer noch sehr kostspielig.'\}} \\
\texttt{VORSCHLAG: \{'C76', 'C56', 'A07', 'C62', 'C81', 'C10', 'C00', 'C08', 'C98'\}} \\
            }
        }
    }
     \\ \\

    \theutterance \stepcounter{utterance}  
    & & \multicolumn{4}{p{0.6\linewidth}}{
        \cellcolor[rgb]{0.9,0.9,0.9}{
            \makecell[{{p{\linewidth}}}]{
                \texttt{\tiny{[P1$\langle$GM]}}
                \texttt{ARGUMENT: \{'C08 hat einen moderaten Aufwand (226) und bietet mir einen guten Wert, was diesen Vorschlag akzeptabler macht als deine vorherigen. Allerdings ist B38 mit 698 Aufwand immer noch sehr kostspielig.'\}} \\
\texttt{VORSCHLAG: \{'C76', 'C56', 'A07', 'C62', 'C81', 'C10', 'C00', 'C08', 'C98'\}} \\
            }
        }
    }
    & \\ \\

    \theutterance \stepcounter{utterance}  
    & \multicolumn{4}{p{0.6\linewidth}}{
        \cellcolor[rgb]{0.9,0.9,0.9}{
            \makecell[{{p{\linewidth}}}]{
                \texttt{\tiny{[P1$\rangle$GM]}}
                \texttt{ARGUMENT: \{'This proposal is getting closer to something workable, but C98 still offers me poor value (262/631). I'd prefer to include C32 instead, which gives me much better value (808/882) despite the higher cost.'\}} \\
\texttt{VORSCHLAG: \{'C76', 'C56', 'A07', 'C62', 'C81', 'C10', 'C00', 'C08', 'C32'\}} \\
            }
        }
    }
    & & \\ \\

    \theutterance \stepcounter{utterance}  
    & & \multicolumn{4}{p{0.6\linewidth}}{
        \cellcolor[rgb]{0.9,0.9,0.9}{
            \makecell[{{p{\linewidth}}}]{
                \texttt{\tiny{[GM$\rangle$P2]}}
                \texttt{ARGUMENT: \{'This proposal is getting closer to something workable, but C98 still offers me poor value (262/631). I'd prefer to include C32 instead, which gives me much better value (808/882) despite the higher cost.'\}} \\
\texttt{VORSCHLAG: \{'C76', 'C56', 'A07', 'C62', 'C81', 'C10', 'C00', 'C08', 'C32'\}} \\
            }
        }
    }
    & \\ \\

    \theutterance \stepcounter{utterance}  
    & & & \multicolumn{4}{p{0.6\linewidth}}{
        \cellcolor[rgb]{0.9,0.9,0.9}{
            \makecell[{{p{\linewidth}}}]{
                \texttt{\tiny{[GM$\langle$P2]}}
                \texttt{ARGUMENT: \{'C32 hat einen extrem hohen Aufwand (882) und würde fast unser gesamtes Budget verbrauchen. Mein letzter Vorschlag mit C98 ist budgetfreundlicher und lässt Raum für andere wertvolle Gegenstände.'\}} \\
\texttt{ABLEHNUNG: \{'C76', 'C56', 'A07', 'C62', 'C81', 'C10', 'C00', 'C08', 'C32'\}} \\
\texttt{VORSCHLAG: \{'C76', 'C56', 'A07', 'C62', 'C81', 'C10', 'C00', 'C08', 'C03'\}} \\
            }
        }
    }
     \\ \\

    \theutterance \stepcounter{utterance}  
    & & \multicolumn{4}{p{0.6\linewidth}}{
        \cellcolor[rgb]{0.9,0.9,0.9}{
            \makecell[{{p{\linewidth}}}]{
                \texttt{\tiny{[P1$\langle$GM]}}
                \texttt{ARGUMENT: \{'C32 hat einen extrem hohen Aufwand (882) und würde fast unser gesamtes Budget verbrauchen. Mein letzter Vorschlag mit C98 ist budgetfreundlicher und lässt Raum für andere wertvolle Gegenstände.'\}} \\
\texttt{ABLEHNUNG: \{'C76', 'C56', 'A07', 'C62', 'C81', 'C10', 'C00', 'C08', 'C32'\}} \\
\texttt{VORSCHLAG: \{'C76', 'C56', 'A07', 'C62', 'C81', 'C10', 'C00', 'C08', 'C03'\}} \\
            }
        }
    }
    & \\ \\

    \theutterance \stepcounter{utterance}  
    & \multicolumn{4}{p{0.6\linewidth}}{
        \cellcolor[rgb]{0.9,0.9,0.9}{
            \makecell[{{p{\linewidth}}}]{
                \texttt{\tiny{[P1$\rangle$GM]}}
                \texttt{ARGUMENT: \{'C03 has mediocre value for me (389/641) while using significant budget. I'll accept C98 instead, which has similar value{-}to{-}cost ratio for me, if we can also include B38 which provides excellent value for its cost.'\}} \\
\texttt{VORSCHLAG: \{'C76', 'C56', 'A07', 'C62', 'C81', 'C10', 'C00', 'C08', 'B38', 'C98'\}} \\
            }
        }
    }
    & & \\ \\

    \theutterance \stepcounter{utterance}  
    & & \multicolumn{4}{p{0.6\linewidth}}{
        \cellcolor[rgb]{0.9,0.9,0.9}{
            \makecell[{{p{\linewidth}}}]{
                \texttt{\tiny{[GM$\rangle$P2]}}
                \texttt{ARGUMENT: \{'C03 has mediocre value for me (389/641) while using significant budget. I'll accept C98 instead, which has similar value{-}to{-}cost ratio for me, if we can also include B38 which provides excellent value for its cost.'\}} \\
\texttt{VORSCHLAG: \{'C76', 'C56', 'A07', 'C62', 'C81', 'C10', 'C00', 'C08', 'B38', 'C98'\}} \\
            }
        }
    }
    & \\ \\

    \theutterance \stepcounter{utterance}  
    & & & \multicolumn{4}{p{0.6\linewidth}}{
        \cellcolor[rgb]{0.9,0.9,0.9}{
            \makecell[{{p{\linewidth}}}]{
                \texttt{\tiny{[GM$\langle$P2]}}
                \texttt{ARGUMENT: \{'Dein Vorschlag überschreitet das Budget: Gesamtaufwand 4926 (max 3629). B38 (698) und C98 (631) zusammen sind zu teuer. Ich schlage vor, stattdessen meinen effizienteren Vorschlag mit C03 zu nehmen.'\}} \\
\texttt{ABLEHNUNG: \{'C76', 'C56', 'A07', 'C62', 'C81', 'C10', 'C00', 'C08', 'B38', 'C98'\}} \\
\texttt{VORSCHLAG: \{'C76', 'C56', 'A07', 'C62', 'C81', 'C10', 'C00', 'C08', 'C03'\}} \\
            }
        }
    }
     \\ \\

    \theutterance \stepcounter{utterance}  
    & & \multicolumn{4}{p{0.6\linewidth}}{
        \cellcolor[rgb]{0.9,0.9,0.9}{
            \makecell[{{p{\linewidth}}}]{
                \texttt{\tiny{[P1$\langle$GM]}}
                \texttt{ARGUMENT: \{'Dein Vorschlag überschreitet das Budget: Gesamtaufwand 4926 (max 3629). B38 (698) und C98 (631) zusammen sind zu teuer. Ich schlage vor, stattdessen meinen effizienteren Vorschlag mit C03 zu nehmen.'\}} \\
\texttt{ABLEHNUNG: \{'C76', 'C56', 'A07', 'C62', 'C81', 'C10', 'C00', 'C08', 'B38', 'C98'\}} \\
\texttt{VORSCHLAG: \{'C76', 'C56', 'A07', 'C62', 'C81', 'C10', 'C00', 'C08', 'C03'\}} \\
            }
        }
    }
    & \\ \\

    \theutterance \stepcounter{utterance}  
    & \multicolumn{4}{p{0.6\linewidth}}{
        \cellcolor[rgb]{0.9,0.9,0.9}{
            \makecell[{{p{\linewidth}}}]{
                \texttt{\tiny{[P1$\rangle$GM]}}
                \texttt{ARGUMENT: \{'You're right about the budget miscalculation. However, C03 provides me significantly less value than other options. I'm willing to drop C00 (780/761) to accommodate B38 (668/698), which gives me comparable value with better budget efficiency.'\}} \\
\texttt{VORSCHLAG: \{'C76', 'C56', 'A07', 'C62', 'C81', 'C10', 'C08', 'B38', 'C03'\}} \\
            }
        }
    }
    & & \\ \\

    \theutterance \stepcounter{utterance}  
    & & \multicolumn{4}{p{0.6\linewidth}}{
        \cellcolor[rgb]{0.9,0.9,0.9}{
            \makecell[{{p{\linewidth}}}]{
                \texttt{\tiny{[GM$\rangle$P2]}}
                \texttt{ARGUMENT: \{'You're right about the budget miscalculation. However, C03 provides me significantly less value than other options. I'm willing to drop C00 (780/761) to accommodate B38 (668/698), which gives me comparable value with better budget efficiency.'\}} \\
\texttt{VORSCHLAG: \{'C76', 'C56', 'A07', 'C62', 'C81', 'C10', 'C08', 'B38', 'C03'\}} \\
            }
        }
    }
    & \\ \\

    \theutterance \stepcounter{utterance}  
    & & & \multicolumn{4}{p{0.6\linewidth}}{
        \cellcolor[rgb]{0.9,0.9,0.9}{
            \makecell[{{p{\linewidth}}}]{
                \texttt{\tiny{[GM$\langle$P2]}}
                \texttt{ARGUMENT: \{'Durch das Entfernen von C00 sparst du zwar Aufwand, aber C00 ist einer meiner wertvollsten Gegenstände (389 Wert). Dein Vorschlag enthält zu viele Gegenstände mit hohem Aufwand. Mein Vorschlag bietet ein besseres Gleichgewicht.'\}} \\
\texttt{ABLEHNUNG: \{'C76', 'C56', 'A07', 'C62', 'C81', 'C10', 'C08', 'B38', 'C03'\}} \\
\texttt{VORSCHLAG: \{'C76', 'C56', 'A07', 'C62', 'C81', 'C10', 'C00', 'C08', 'C98'\}} \\
            }
        }
    }
     \\ \\

    \theutterance \stepcounter{utterance}  
    & & \multicolumn{4}{p{0.6\linewidth}}{
        \cellcolor[rgb]{0.9,0.9,0.9}{
            \makecell[{{p{\linewidth}}}]{
                \texttt{\tiny{[P1$\langle$GM]}}
                \texttt{ARGUMENT: \{'Durch das Entfernen von C00 sparst du zwar Aufwand, aber C00 ist einer meiner wertvollsten Gegenstände (389 Wert). Dein Vorschlag enthält zu viele Gegenstände mit hohem Aufwand. Mein Vorschlag bietet ein besseres Gleichgewicht.'\}} \\
\texttt{ABLEHNUNG: \{'C76', 'C56', 'A07', 'C62', 'C81', 'C10', 'C08', 'B38', 'C03'\}} \\
\texttt{VORSCHLAG: \{'C76', 'C56', 'A07', 'C62', 'C81', 'C10', 'C00', 'C08', 'C98'\}} \\
            }
        }
    }
    & \\ \\

    \theutterance \stepcounter{utterance}  
    & \multicolumn{4}{p{0.6\linewidth}}{
        \cellcolor[rgb]{0.9,0.9,0.9}{
            \makecell[{{p{\linewidth}}}]{
                \texttt{\tiny{[P1$\rangle$GM]}}
                \texttt{ARGUMENT: \{'I understand C00 is valuable to you, but C98 (262/631) is very inefficient for me. I'll accept your proposal if we replace C98 with C32 (808/882), which gives me much better value despite the higher cost.'\}} \\
\texttt{VORSCHLAG: \{'C76', 'C56', 'A07', 'C62', 'C81', 'C10', 'C00', 'C08', 'C32'\}} \\
            }
        }
    }
    & & \\ \\

    \theutterance \stepcounter{utterance}  
    & & \multicolumn{4}{p{0.6\linewidth}}{
        \cellcolor[rgb]{0.9,0.9,0.9}{
            \makecell[{{p{\linewidth}}}]{
                \texttt{\tiny{[GM$\rangle$P2]}}
                \texttt{ARGUMENT: \{'I understand C00 is valuable to you, but C98 (262/631) is very inefficient for me. I'll accept your proposal if we replace C98 with C32 (808/882), which gives me much better value despite the higher cost.'\}} \\
\texttt{VORSCHLAG: \{'C76', 'C56', 'A07', 'C62', 'C81', 'C10', 'C00', 'C08', 'C32'\}} \\
            }
        }
    }
    & \\ \\

    \theutterance \stepcounter{utterance}  
    & & & \multicolumn{4}{p{0.6\linewidth}}{
        \cellcolor[rgb]{0.9,0.9,0.9}{
            \makecell[{{p{\linewidth}}}]{
                \texttt{\tiny{[GM$\langle$P2]}}
                \texttt{ARGUMENT: \{'C32 hat mit 882 einen zu hohen Aufwand und würde die Budgetgrenze sprengen. Selbst ohne C98 würde C32 allein den Gesamtaufwand auf 3654 bringen, was über dem Maximum von 3629 liegt.'\}} \\
\texttt{ABLEHNUNG: \{'C76', 'C56', 'A07', 'C62', 'C81', 'C10', 'C00', 'C08', 'C32'\}} \\
\texttt{VORSCHLAG: \{'C76', 'C56', 'A07', 'C62', 'C81', 'C10', 'C00', 'C08', 'C03'\}} \\
            }
        }
    }
     \\ \\

    \theutterance \stepcounter{utterance}  
    & & \multicolumn{4}{p{0.6\linewidth}}{
        \cellcolor[rgb]{0.9,0.9,0.9}{
            \makecell[{{p{\linewidth}}}]{
                \texttt{\tiny{[P1$\langle$GM]}}
                \texttt{ARGUMENT: \{'C32 hat mit 882 einen zu hohen Aufwand und würde die Budgetgrenze sprengen. Selbst ohne C98 würde C32 allein den Gesamtaufwand auf 3654 bringen, was über dem Maximum von 3629 liegt.'\}} \\
\texttt{ABLEHNUNG: \{'C76', 'C56', 'A07', 'C62', 'C81', 'C10', 'C00', 'C08', 'C32'\}} \\
\texttt{VORSCHLAG: \{'C76', 'C56', 'A07', 'C62', 'C81', 'C10', 'C00', 'C08', 'C03'\}} \\
            }
        }
    }
    & \\ \\

    \theutterance \stepcounter{utterance}  
    & \multicolumn{4}{p{0.6\linewidth}}{
        \cellcolor[rgb]{0.9,0.9,0.9}{
            \makecell[{{p{\linewidth}}}]{
                \texttt{\tiny{[P1$\rangle$GM]}}
                \texttt{ARGUMENT: \{'I see the budget constraint is very tight. Let me calculate carefully: your proposal totals 3424, which works. While C03 isn't ideal for me, I'm willing to accept this compromise to reach an agreement.'\}} \\
\texttt{ZUSTIMMUNG: \{'C76', 'C56', 'A07', 'C62', 'C81', 'C10', 'C00', 'C08', 'C03'\}} \\
            }
        }
    }
    & & \\ \\

    \theutterance \stepcounter{utterance}  
    & & & \multicolumn{2}{p{0.3\linewidth}}{
        \cellcolor[rgb]{0.9,0.9,0.9}{
            \makecell[{{p{\linewidth}}}]{
                \texttt{\tiny{[GM$|$GM]}}
                \texttt{game successful} \\
            }
        }
    }
    & & \\ \\

    \theutterance \stepcounter{utterance}  
    & & & \multicolumn{2}{p{0.3\linewidth}}{
        \cellcolor[rgb]{0.9,0.9,0.9}{
            \makecell[{{p{\linewidth}}}]{
                \texttt{\tiny{[GM$|$GM]}}
                \texttt{end game} \\
            }
        }
    }
    & & \\ \\

\end{supertabular}
}

\end{document}
